\begin{statement}{D.55}{Proved}{The Legendre-relative prosolvable Fermat assertion and the three-power prosolvable tower}
For an odd integer \(n\ge3\), let
\[
 U_n=\{X^n+Y^n+Z^n=0,\ XYZ\ne0\},\qquad
 N=\pi_1(U_{n,\overline{\mathbb Q}}),\qquad
 \Pi_n=\pi_1(U_n).
\]
For an arithmetic-stable open normal subgroup \(R\subset N\), put
\(R^{(\infty)}=\bigcap_{m\ge1}R^{(m)}\), with closed derived groups.
Then the following statements hold.
\begin{enumerate}
\item If \(\ell\ge5\) is prime and divides \(n\), let
\(R_\ell=\ker(N\to\operatorname{SL}_2(\mathbb F_\ell))\)
for the Legendre family pulled back by \(\lambda=-(X/Z)^n\).
Every section of \(\Pi_n/R_\ell^{(\infty)}\to G_{\mathbb Q}\)
is cuspidal.  In particular the Legendre-relative prosolvable
clause of the Fermat hunch holds at every \(5^r\) level, using
\(\ell=5\).
\item If \(n=3^r\), \(r\ge1\), every section of
\(\Pi_n/R^{(\infty)}\to G_{\mathbb Q}\) is cuspidal for every
such \(R\).  This includes the plain geometrically prosolvable
quotient by taking \(R=N\).
\item Consequently the noncuspidal full metabelian families in
D.20 have no compatible lift to the plain prosolvable inverse
limit on the three-power tower, or to the indicated
Legendre-relative prosolvable inverse limit on the five-power
tower.  No first obstructed finite derived length is asserted.
\end{enumerate}
More generally, over every finite extension of \(\mathbb Q_p\),
D.54 implies the point-or-cusp conclusion for the relative
prosolvable quotient above any fixed arithmetic-stable finite
geometric Galois cover.  The proof below includes the arithmetic
twist which makes that assertion exact.
\end{statement}
\begin{proof}
\paragraph{Step 1: the exact arithmetic twist.}

Let $U/k$ be hyperbolic, $\Pi_U=\pi_1(U)$ and
$N=\pi_1(U_{\bar k})$.  Let $R\subset N$ be open and normal,
stable under $\Pi_U$-conjugation, and put
\[
 R^{(\infty)}=\bigcap_{m\ge1}R^{(m)},\qquad
 \Pi_{U,R}^{\mathrm{rel,sol}}
      =\Pi_U/R^{(\infty)}.
\]
All derived subgroups are closed.  The first intersection is the
kernel of the maximal prosolvable quotient of $R$: every finite
solvable image kills some $R^{(m)}$, and the quotient by the
intersection embeds in the product of the finite-derived quotients.

Fix an actual section $s:G_k\to\Pi_{U,R}^{\mathrm{rel,sol}}$.
The subgroup
\[
             H=R^{\mathrm{sol}}\,s(G_k)
                 \subset\Pi_{U,R}^{\mathrm{rel,sol}}
\]
is open, projects onto $G_k$, and has intersection
$R^{\mathrm{sol}}$ with the geometric subgroup.  Its preimage in
$\Pi_U$ defines a finite \'etale cover $Y_s\to U$ over $k$, with
constant field exactly $k$.  Geometrically it is the Galois cover
corresponding to $R$.  Since its full geometric fundamental group
is $R$, its geometrically prosolvable fundamental group is exactly
$H$.  In particular the section $s$ is an actual prosolvable section
of $Y_s/k$; no extension of $k$ and no descent of a geometric
conjugator is required.

For $k/\mathbb Q_p$ finite, D.54 makes this section
point or cusp induced on $Y_s$.  Its composite in
$\Pi_{U,R}^{\mathrm{rel,sol}}$ is therefore point or cusp induced
on $U$.  This proves the local relative assertion for \emph{every}
fixed finite geometric cover, including nonsolvable ones.

\paragraph{Step 2: finite support using only solvable covers.}

We use the following primary version to avoid any implicit real-place
assumption.  \cite[Theorem 4.1]{DE_Hoshi2012} applies to a locally
geometric pro-$\mathcal C$ section over a number field when
$\mathcal C$ contains all finite solvable groups.  It is geometric
if its underlying local compactification points lie in one finite
closed subscheme at all finite places outside a density-zero set.
\cite[Corollary 4.2]{DE_Hoshi2012} in particular proves geometricity when the curve
has a nonconstant map to an abelian variety $A$ with $A(k)$ and
$\operatorname{Sha}(A/k)$ finite.  Local geometricity in this statement is at
finite places.  The proof passes to a totally imaginary extension,
uses Stoll's Theorem 8.2 at every solvable section neighborhood,
and descends by Hoshi's Lemma 1.6.

Here is the cohomological content of the special case used below.
The induced section on $A$ is a class in
$H^1(k,T A)$, and its finite local restrictions are point Kummer
classes.  Therefore it lies in the inverse-limit finite Selmer
group.  The quotient of the finite-level Selmer group by
$A(k)/nA(k)$ is killed uniformly, using finiteness of $\operatorname{Sha}$ and
the fact that real local groups have exponent at most two.
Taking the inverse limit over all $n$ kills that quotient.  Hence
the class belongs to $A(k)$ if $A(k)$ is finite.  Full local Tate
Kummer injectivity makes each underlying local curve point lie in
the finite fibre above this one global point of $A$.  Stoll's
finite-support theorem \cite[Theorem 8.2]{DE_Stoll2007}, with the
finite-place correction in \cite{DE_StollErratum}, then supplies rational points on every
solvable section neighborhood.  Their compatible inverse limit
gives a point or cusp decomposition group.  Thus full finite
non-solvable torsors are not required in this argument.

For the modular target used below, Stoll's Corollary 8.8 and its
proof \cite{DE_Stoll2007} provide exactly the needed abelian variety: every
positive-genus $X_0(M)/\mathbb Q$ has a nontrivial quotient of its
Jacobian that is modular of analytic rank zero.  Its rational
points and Tate--Shafarevich group are finite.  The map of the
modular curve to that quotient is nonconstant, since the curve's
Abel--Jacobi image generates its Jacobian.  Thus any hyperbolic
curve with a nonconstant map to such an $X_0(M)$ satisfies the
preceding prosolvable finite-support principle.

\paragraph{Step 3: the Legendre-relative Fermat quotient.}

Let $\ell\ge5$ be prime, let $n$ be odd and divisible by $\ell$,
and let
\[
 U_n:\quad X^n+Y^n+Z^n=0,\qquad XYZ\ne0,
 \qquad\lambda=-(X/Z)^n.
\]
Pull the Legendre family $y^2=x(x-1)(x-\lambda)$ back to $U_n$.
Let $N=\pi_1(U_{n,\overline{\mathbb Q}})$, and set
\[
 R_\ell=\ker\bigl(N\longrightarrow
                         \operatorname{SL}_2(\mathbb F_\ell)\bigr).
\]
This geometric image is the whole $\operatorname{SL}_2$.
Indeed the tripod image is generated by its two standard
unipotent matrices and is $\operatorname{SL}_2(\mathbb F_\ell)$;
$N$ is normal in the tripod group with abelian quotient, whereas
$\operatorname{SL}_2(\mathbb F_\ell)$ is perfect for $\ell\ge5$.
The residual kernel is stable under the arithmetic group.

\medskip
\noindent\textbf{The required conclusion.}
Every section of
\[
 \Pi_{U_n}/R_\ell^{(\infty)}\longrightarrow G_{\mathbb Q}
                                                        \tag{D.55.1}
\]
is cuspidal.  In particular this settles the Legendre-relative
part of the Fermat hunch for the $5^r$ tower, using $\ell=5$.

To prove it, take a section $s$ of (D.55.1).  Step 1 makes every
finite local restriction point or cusp induced.  Consequently
every local residual calculation in D.23 applies to its actual
Legendre representation $\bar\rho_s$.  For clarity the divisibility
used there is $\ell\mid n$:
\[
  v_p(j)=8v_p(2)-2n|v_p(X/Z)|,
       \qquad [q(s)]_\ell=[256^{-1}]_\ell
\]
at a cusp, with the analogous formula using $Y/Z$ when necessary.
After one global quadratic twist, the representation is unramified
outside $\{2,\ell\}$, has conductor at most two away from $\ell$,
is finite flat of Serre weight two at $\ell$, and has determinant
$\bar\chi_\ell$.  It is odd because complex conjugation has
determinant $-1$.  An irreducible representation would therefore
come from a weight-two cusp form of level one or two, whose spaces
vanish.  The representation is reducible, exactly as in the
independently audited D.23 proof.  This argument uses the full
cuspidal constant $256^{-1}$; it does not presume that the cusp
representation splits.

Its invariant line gives a section-supported component $Z$ of
the finite cover of $U_n$ parametrizing a cyclic subgroup of order
$\ell$ of the Legendre curve.  Its geometric fundamental group
$R_B\subset N$ is the inverse image of a Borel subgroup
$B\subset\operatorname{SL}_2(\mathbb F_\ell)$, up to conjugacy.
Since $B$ is solvable,
\[
                       R_B^{(\infty)}=R_\ell^{(\infty)}.    \tag{D.55.2}
\]
For the two inclusions, any solvable quotient of $R_B$ restricts
to a solvable quotient of $R_\ell$, while
$R_B/R_\ell^{(\infty)}$ is an extension of the prosolvable group
$R_\ell^{\mathrm{sol}}$ by the finite solvable group $B$ and is
itself prosolvable.  Thus the subgroup of (D.55.1) belonging to $Z$
is \emph{exactly} $\Pi_Z^{\mathrm{sol}}$.  The lifted section is
a prosolvable section of this actual curve, not just a section of
an unspecified relative quotient.

As in D.23, $Z$ has a nonconstant map to $Y_0(4\ell)$: quotient
the Legendre curve by one rational two-torsion point and combine
the image of its order-$\ell$ subgroup with the cyclic order-four
subgroup supplied by the other rational two-torsion point.
The map extends to the compactifications, and
$g(X_0(4\ell))=(\ell-3)/2>0$.  D.54 makes the
section of $Z$ locally geometric.  Step 2, with the rank-zero
modular abelian quotient, makes it globally geometric.  Its image
in (D.55.1) is therefore a point or cusp section of $U_n$.
The prime-exponent argument of D.23, or the coordinate power map
$U_n\to U_\ell$ followed by that argument, gives $U_n(\mathbb Q)=
\varnothing$.  Thus the section is cuspidal, proving the stated
conclusion.

\paragraph{Step 4: the three-power tower and the exact limit limitation.}

For $n=3^r$, the compactification of $U_n$ has a nonconstant map
to the Fermat cubic.  The explicit cubic transformation in D.23
identifies that cubic with $X_0(27)$, whose rational points and
Tate--Shafarevich group are finite.  Therefore Step 2 and the
local candidate already show that every section of
$\Pi_{U_n}^{\mathrm{sol}}$ is cuspidal.  The same is true of
$\Pi_{U_n}/R^{(\infty)}$ for \emph{any} arithmetic-stable finite
index geometric subgroup $R$: its arithmetic twist $Y_s$ from
Step 1 still has a nonconstant map to $X_0(27)$, so its
prosolvable section is geometric.  The absence of ordinary
$U_n(\mathbb Q)$ points was checked in D.23.

Consequently the noncuspidal full metabelian classes of D.20 on
the three-power tower have no compatible lift to the full
prosolvable inverse limit.  For the five-power tower they have no
compatible lift to the relative prosolvable quotient (D.55.1).
Neither assertion identifies a first finite derived length.
In particular, over $G_{\mathbb Q}$ the sets of continuous sections
need not be compact when ramification is unrestricted; existence
of an independent lift at each finite derived length cannot be
silently replaced by existence of one compatible family.

The phrase ``only cuspidal'' here concerns the Fermat curves.
It does not assert this for the tripod itself, which has ordinary
rational point sections.  Also, the proof of Step 3 requires
$\ell\mid n$ and does not assert the same conclusion for an
unrelated Legendre prime without an additional argument.

\end{proof}
